 % 本文件由 scripts/assemble.py 自动装配，勿手改（改 parts/ 后重跑）
% 第11章 超导电性

% 第11章 SUPERCONDUCTIVITY
% 页码核对说明：经与扫描页核对，第 11 章实际始于书页 379（p393.png）；
% p389–p392（书页 375–378）为第 10 章 §10.12 的结尾，归第 10 章文件。
% 本文件实际覆盖书页 379–385（p393–p399）：章首 + §11.1 + §11.2。

\chapter{超导电性}

\epigraphbook{既然变化本是自然的法则，\\ 唯独恒常反而显得奇怪。}{EARL OF ROCHESTER}

\section{电子间的吸引}

超导电性（superconductivity）长期以来被看作金属诸性质中最异乎寻常、最神秘的一种；但是 Bardeen、Cooper 与 Schrieffer 的理论——即\emph{BCS 理论}——已经解释了如此之多，以致我们可以说，如今我们对超导态的理解，几乎不亚于对“正常”态的理解。我们并不打算在这里论及这个庞大而进展迅速的领域的全部内容；重点将放在那些原子过程上，许多不寻常的宏观现象正是由它们引发的。

\begin{figure}[H]
\centering
\includegraphics[width=0.72\textwidth]{fig_199.pdf}
\caption*{图 199\quad 电子附近晶格的极化。}
\end{figure}

整个效应源自任意两个能量几乎相同的电子之间的一种微弱吸引力。我们通常认为自由电子通过其库仑相互作用而相互\emph{排斥}；不过，正如 §5.8 所示，在长距离上这个场会因“其他”电子的屏蔽（screening）效应而显著减弱。但在晶格中，电子倾向于把正离子拉向自己，从而被一个晶格比平常略为密集的区域所包围。进入其邻近区域的另一个电子将被引向这个区域；看上去就好像它被第一个电子吸引着。打个比方：两个粒子紧挨着一同坐进床垫的同一个凹坑里，是能够占到便宜的。

计算这种力，最简单的办法是把它描述为一个电子发射“虚”声子、另一个电子吸收此声子的过程。考察图 200(a) 所定义的过程：处于状态 $\mathbf{k}$ 的电子发射一个声子，被散射到状态 $\mathbf{k}-\mathbf{K}$；处于状态 $\mathbf{k}'$ 的电子吸收这个声子，被散射到 $\mathbf{k}'+\mathbf{K}$。

\begin{figure}[H]
\centering
\includegraphics[width=0.72\textwidth]{fig_200.pdf}
\caption*{图 200\quad 通过声子交换的电子–电子相互作用。}
\end{figure}

在二级微扰论中，我们应当用矩阵元
\begin{equation*}
\langle \mathbf{k}-\mathbf{K},\,\mathbf{k}'+\mathbf{K}\,|\,V\,|\,\mathbf{k},\,\mathbf{k}'\rangle
=\frac{\mathcal{M}_{\mathbf{k},\,\mathbf{k}-\mathbf{K}}\,\mathcal{M}^{*}_{\mathbf{k}',\,\mathbf{k}'+\mathbf{K}}}{\mathcal{E}(\mathbf{k})-\mathcal{E}(\mathbf{k}-\mathbf{K})-\hbar\nu_{\mathbf{q}}}.
\tag{11.1}
\end{equation*}
来表示这一跃迁。在这个表达式中，$\mathcal{M}_{\mathbf{k},\mathbf{k}-\mathbf{K}}$ 代表电子–声子相互作用的矩阵元，与 §6.13 中的相同。能量分母代表由初态到中间态的能量变化，其中包括所产生的声子的能量 $\hbar\nu_{\mathbf{q}}$；对 N 过程有 $\mathbf{q}=\mathbf{K}$，不过电子–声子 U 过程也可以包括在内。

还有另一个过程对总散射作出贡献，如图 200(b) 所示：波矢为 $-\mathbf{q}$ 的声子先由 $\mathbf{k}'$ 发射出来，然后被 $\mathbf{k}$ 吸收。它给出一个与式 (11.1) 类似的表达式，只是分母稍有不同。把这两个过程合并起来，并假定电子–声子矩阵元实质上相同，我们就得到
\begin{equation*}
\langle \mathbf{k}-\mathbf{K},\,\mathbf{k}'+\mathbf{K}\,|\,V\,|\,\mathbf{k},\,\mathbf{k}'\rangle
=\frac{|\mathcal{M}_{\mathbf{k},\,\mathbf{k}-\mathbf{K}}|^{2}\,2\hbar\nu_{\mathbf{q}}}{\{\mathcal{E}(\mathbf{k})-\mathcal{E}(\mathbf{k}-\mathbf{K})\}^{2}-(\hbar\nu_{\mathbf{q}})^{2}}.
\tag{11.2}
\end{equation*}
它对电子所产生的效果，就好像两者之间存在一种直接的相互作用，其 Fourier 系数 $V(\mathbf{K})$ 由式 (11.2) 给出。它通常是正的；但对于
\begin{equation*}
|\mathcal{E}(\mathbf{k})-\mathcal{E}(\mathbf{k}-\mathbf{K})|<\hbar\nu_{\mathbf{q}}
\tag{11.3}
\end{equation*}
这一狭窄的能量范围，相互作用是负的，相应于一种吸引力。在式 (11.3) 的范围内，这个力具有 Fourier 分量
\begin{equation*}
V(\mathbf{K})\approx-\frac{2\,|\mathcal{M}_{\mathbf{k},\,\mathbf{k}-\mathbf{K}}|^{2}}{\hbar\nu_{\mathbf{q}}}.
\tag{11.4}
\end{equation*}
这或许算不上一种很强的力。我们是在极低的温度下工作，因此可以把声子发射过程看作由晶格的零点运动所诱发。利用式 (2.109)、(6.86) 与 (6.93)，我们得到
\begin{align*}
|\mathcal{M}_{\mathbf{k},\,\mathbf{k}-\mathbf{K}}|^{2}
&=|\mathcal{U}(\mathbf{K})|^{2}\,|\mathbf{K}\cdot\mathbf{U}_{\mathbf{q}}|^{2}\\
&=|\mathcal{U}(\mathbf{K})|^{2}\,\frac{\mathbf{K}^{2}\hbar}{2NM\nu_{q}},
\tag{11.5}
\end{align*}
式中 $\mathcal{U}(\mathbf{K})$ 是单个离子的屏蔽势的 Fourier 变换。如果只有 N 过程，并且声速由式 (6.84) 正确给出，那么
\begin{align*}
V(\mathbf{K}) &\approx-\frac{|\mathcal{U}(\mathbf{K})|^{2}\,\mathbf{q}^{2}}{NM\nu_{q}^{2}}\\
&=-\frac{|\mathcal{U}(\mathbf{K})|^{2}}{D_{0}s^{2}}=-\frac{3}{2}\,\frac{|\mathcal{U}(\mathbf{K})|^{2}}{NZ\mathcal{E}_{F}}.
\tag{11.6}
\end{align*}
这样，在 N 过程所允许的 $\mathbf{K}$ 值范围内，有效相互作用并不迅速减小。对更大的 $K$ 值，U 过程会使 $V(\mathbf{K})$ 有所增大，不过，正如 §6.13 和 §7.5 中所讨论的，电子 Bloch 态中平面波的混合避免了 $\mathbf{K}=\mathbf{g}$ 处的灾难。显然，这种吸引是短程的——直到两电子之间所能交换动量的最大值，它都具有很大的矩阵元。

当然，电子之间还存在着普通的库仑排斥。如同式 (5.32) 那样，相应的矩阵元应为
\begin{equation*}
V_{\mathrm{Coul.}}(\mathbf{K})=\frac{4\pi e^{2}}{K^{2}+\lambda^{2}},
\tag{11.7}
\end{equation*}
式中 $\lambda$ 是式 (5.22) 的屏蔽参量，或者是像式 (5.36) 那样的某种更复杂的表达式。我们把式 (11.4) 与式 (11.7) 之和取作对 \emph{Fröhlich 相互作用}（Fröhlich interaction）的估计。超导电性的判据就是它应为负值。用电子–声子相互作用来表述，这个判据并没有任何很简单的形式；不过有证据表明，Bardeen 公式 (6.93) 中“赝势部分”$\mathcal{V}'(K)$ 的强度才是关键因素。

上面的论证已经足够简单；但要想恰当地为它辩护，就需要考察电子–声子联合系统的能量中的各种贡献，其中包括 §6.11 中声速的重整化。应当注意，我们处理的只是像式 (11.3) 中那样处于很窄相对能量范围内的电子；在这个范围之外，晶格的运动太慢，跟不上电子的运动。但由于对净相互作用最重要的贡献来自相当大的 $K$ 值——相应于大的 $q$ 值——这个范围的量级为 $k\Theta$，这里 $\Theta$ 是 Debye 温度。既然超导电性只在极低温度下出现，这个能量范围就远比费米面（Fermi surface）上“热层”的厚度 $kT$ 宽广得多。因此，对于那些真正能够参与超导转变的电子来说，近似式 (11.4) 是合理的。

\section{Cooper 对}

为了领会电子之间交换声子所产生的吸引力的效果，我们现在来计算总散射振幅 $t$，即这样一个矩阵元：它的平方给出下述过程的概率——处于 $\mathbf{k}_1$ 和 $\mathbf{k}_2$ 两态的两个电子被散射到另外两个态 $\mathbf{k}_1'$ 和 $\mathbf{k}_2'$。如图 201(a) 那样，我们将假定这些态正好位于费米球（Fermi sphere）之外。

对散射最显而易见的贡献将是式 (11.2) 所定义的直接相互作用 $V(\mathbf{K})$。它对我们眼下的计算来说未免太复杂，所以让我们抽出其中的吸引部分，并假定它是一个常数。于是，我们的一级散射振幅将取为
\begin{equation*}
t_{1}=V(\mathbf{K})=V,
\tag{11.8}
\end{equation*}
其中 $V$ 是一个负的常数，与 $\mathbf{k}_1'$ 和 $\mathbf{k}_1$ 无关。但这个吸引势只对能量几乎相同的电子才成立；为了把条件 (11.3) 纳入理论，我们将约定：除非
\begin{equation*}
|\mathcal{E}(\mathbf{k}_{1})-\mathcal{E}(\mathbf{k}_{2})|<w,
\tag{11.9}
\end{equation*}
否则 $V(\mathbf{K})$ 为零，这里 $w$ 是一个常数。由式 (11.3) 显然可见，$w$ 应当具有 $k\Theta$ 的量级。

通常在散射问题中，近似式 (11.8) 已经足够。然而，在 $t$ 的微扰展开中还会有一些进一步的项，对应于电子从 $\mathbf{k}_1$ 出发、终于 $\mathbf{k}_1'$ 的逐次虚过程。例如，下一项可能就是图 201(b) 所示的过程：我们在初态与终态之间插入一个虚态，其中 $\mathbf{k}_1$ 先变到 $\mathbf{k}_1+\mathbf{K}$、$\mathbf{k}_2$ 先变到 $\mathbf{k}_2-\mathbf{K}$，然后再作第二次跃迁，变到 $\mathbf{k}_1'$ 和 $\mathbf{k}_2'$。

\begin{figure}[H]
\centering
\includegraphics[width=0.72\textwidth]{fig_201.pdf}
\caption*{图 201\quad (a) 电子的直接散射。(b) 二阶项。}
\end{figure}

这种过程对散射振幅的贡献由通常的微扰表达式给出：
\begin{equation*}
t_{2}(\mathbf{K})=\frac{V^{2}}{\mathcal{E}(\mathbf{k}_{1})+\mathcal{E}(\mathbf{k}_{2})-\mathcal{E}(\mathbf{k}_{1}+\mathbf{K})-\mathcal{E}(\mathbf{k}_{2}-\mathbf{K})}.
\tag{11.10}
\end{equation*}
但这一项只在非常有限的条件下才存在。由于 Pauli 原理，中间态不得位于费米面以下；而且按式 (11.9)，每次跃迁的能量差必须小于 $w$。于是我们得到如下规则：
\begin{equation*}
\mathcal{E}_{F}<\mathcal{E}(\mathbf{k}_{1}+\mathbf{K})<\mathcal{E}(\mathbf{k}_{1})+w,\quad
\mathcal{E}_{F}<\mathcal{E}(\mathbf{k}_{2}-\mathbf{K})<\mathcal{E}(\mathbf{k}_{2})+w,
\tag{11.11}
\end{equation*}
等等；如果式 (11.10) 还想对 $t_2$ 有任何贡献的话，这些规则就必须得到满足。这些条件类似于 Kondo 效应（§10.6）中虚自旋翻转跃迁的选择定则；在同时计及直接过程和交换过程时，各个态的占据数并\emph{不}会从二阶项中消去。

这些规则限制性如此之强，以致通常我们可以忽略 $t_2$。但当 $\mathbf{k}_1$ 和 $\mathbf{k}_2$ 的方向近乎相反时，条件 (11.11) 的两式合二为一，对可能发生散射的 $\mathbf{K}$ 值范围也就不再施加那么苛刻的限制。事实上，此时 $t_2$ 可以变得相当大。

为说明这一点，方便的做法是从费米能级（Fermi level）起量度能量。我们对每个态定义
\begin{align*}
\xi(\mathbf{k}) &\equiv \mathcal{E}(\mathbf{k})-\mathcal{E}_{F}\\
&\approx \hbar v_{F}(k-k_{F})
\tag{11.12}
\end{align*}
（只要 $k$ 近似等于 $k_F$）。我们还引入两个具有能量量纲的新辅助变量
\begin{equation*}
\Delta=\xi(\mathbf{k}_{1})+\xi(\mathbf{k}_{2}),\qquad
\eta=\hbar v_{F}\,|\mathbf{k}_{1}+\mathbf{k}_{2}|
\tag{11.13}
\end{equation*}
它们量度的是我们这对电子的总能量和总动量。

为得到二阶微扰项，我们把式 (11.10) 在满足式 (11.11) 的所有 $\mathbf{K}$ 值上积分。利用式 (11.12) 和 (11.13)，这可以化成一个更简单的公式：
\begin{align*}
t_{2} &= V^{2}\int \frac{\mathrm{d}\mathbf{K}}{\mathcal{E}(\mathbf{k}_{1})+\mathcal{E}(\mathbf{k}_{2})-\mathcal{E}(\mathbf{k}_{1}+\mathbf{K})-\mathcal{E}(\mathbf{k}_{2}-\mathbf{K})}\\
&= \frac{2\pi V^{2}k_{F}^{2}}{8\pi^{3}\hbar v_{F}}\int_{0}^{w}\mathrm{d}\xi\int_{-1}^{1}\frac{\mathrm{d}(\cos\theta)}{\Delta+\eta\cos\theta-2\xi}.
\tag{11.14}
\end{align*}

\begin{figure}[H]
\centering
\includegraphics[width=0.72\textwidth]{fig_202.pdf}
\caption*{图 202\quad (a) 在 $\mathbf{k}_3$ 处产生一个空穴、随后又被 $\mathbf{k}_2$ 填充的过程。(b) 逐次的对散射。(c) 涉及空穴的对散射。}
\end{figure}

当 $\Delta$ 和 $\eta$ 都远小于 $w$ 时，被积函数中出现一个奇点，它使积分呈对数行为（可参照式 (10.56)）：
\begin{equation*}
t_{2}\approx-\frac{V^{2}}{4\pi^{2}}\,\frac{k_{F}^{2}}{\hbar v_{F}}\ln\frac{2w}{\max{(\Delta,\eta)}}.
\tag{11.15}
\end{equation*}
通常这一项远小于 $t_1$，原本可以忽略。但如果 $\Delta$ 和 $\eta$ 恰好都非常小，式 (11.15) 就可能变得很大——这时我们或许不得不考虑 $t$ 的微扰级数中更进一步的项。

这样的项可能非常复杂，因为我们可以考虑如图 202(a) 那样的虚过程：一个来自费米面内部的电子与这对电子中的一个发生相互作用，留下一个空穴，而这个空穴最终在此后的某个过程中被填满。不过，这类过程在微扰级数中给出的贡献往往小得多。重要的是那些把图 201(b) 所示过程一再重复的情形——电子对不断地自我散射，如图 202(b) 那样。每一次这样的虚散射都引入一个类似于式 (11.15) 的因子，当 $\Delta$ 和 $\eta$ 很小时还包括那个大的对数积分。

为了得到正确的结果，我们还必须把这样的过程包括进来：一对电子从费米海（Fermi sea）中被散射出去，留下一对空穴，它们可以接收 $\mathbf{k}_1$ 和 $\mathbf{k}_2$ 处的电子（图 202(c)）。这一过程对 $t_2$ 给出与式 (11.15) 相同的贡献。把所有这些项加起来，我们看到散射振幅必定表现为无穷级数
\begin{equation*}
t\approx V+\gamma V^{2}+\gamma^{2}V^{3}+\ldots,
\tag{11.16}
\end{equation*}
其中
\begin{equation*}
\gamma=-\frac{1}{2\pi^{2}}\,\frac{k_{F}^{2}}{\hbar v_{F}}\ln\frac{2w}{\max{(\Delta,\eta)}}.
\tag{11.17}
\end{equation*}
但级数 (11.16) 可以求和：
\begin{equation*}
t\approx\frac{V}{1-\gamma V}.
\tag{11.18}
\end{equation*}
这样，散射振幅的行为就好像它含有一个极点。这是众所周知的一个征兆，表明自由粒子彼此散射的假设已经失效；它正是两个电子形成\emph{束缚态}（bound state）的标志。假设我们取 $\eta=0$，使得 $\mathbf{k}_1=-\mathbf{k}_2$。极点出现时的能量由下式确定：
\begin{align*}
1&=\gamma V\\
&=-\frac{V}{2\pi^{2}}\,\frac{k_{F}^{2}}{\hbar v_{F}}\ln\frac{2w}{\Delta}.
\tag{11.19}
\end{align*}
（译注：原书式 (11.19) 第二行分母印作 $v_F$，漏排 $\hbar$；由式 (11.17) 显见应有 $\hbar v_F$，译文已补正。）

换句话说，当这对电子的总能量在费米能级以下等于
\begin{equation*}
\Delta_{0}=2w\exp\left(-\frac{2\pi^{2}\hbar v_{F}}{|V|\,k_{F}^{2}}\right)
\tag{11.20}
\end{equation*}
时，它们就倾向于彼此束缚而形成一个“准分子”（quasi-molecule）。这样的状态称为 \emph{Cooper 对}（Cooper pair）。

这类态有可能存在，但这并没有解决超导电性的问题。我们可以这样论证：金属中所有的电子都倾向于结成这样的对——$\mathbf{k}$ 与 $-\mathbf{k}$ 相配，而且这种由 Bose–Einstein 型“准分子”组成的气体随后便能发生 Bose–Einstein 凝聚。然而这些电子对并不稳定；作为对超导态的描述，这种说法有启发性，却缺乏说服力。

最重要的结果是关于电子对能量的公式 (11.20)。我们可以采用自由电子模型，它给出
\begin{align*}
\frac{k_{F}^{2}}{2\pi^{2}\hbar v_{F}} &= \frac{1}{2}\,\mathcal{N}(\mathcal{E}_{F})\\
&= \frac{3}{4}\,\frac{ZN}{\mathcal{E}_{F}}
\tag{11.21}
\end{align*}
（在 BCS 理论中，这个量记作 $N(0)$，即费米能级处\emph{每个自旋}的态密度）。把式 (11.6) 代入式 (11.20)，并取 $w\approx k\Theta$，就应当得到
\begin{align*}
\Delta_{0} &\approx 2k\Theta\exp\{-2/\mathcal{N}(\mathcal{E}_{F})\,|V|\}\\
&\approx 2k\Theta\exp\{-8\mathcal{E}_{F}^{2}/9\,|\mathcal{U}|^{2}\}.
\tag{11.22}
\end{align*}
如果我们假定 $\Delta_{0}$ 是超导转变能量的某种量度，并令它等于 $kT_{c}$，就会发现转变温度 $T_{c}$ 应当介于 Debye 温度的 1/10 与 1/10,000 之间，视电子–离子相互作用势 $|\mathcal{U}|$ 的大小而定。实际观察到的正是如此——我们现在可以理解，为什么 $T_{c}$ 与 $\Theta$ 相差得如此悬殊。

由这个论证还可得出的另一个效应是\emph{同位素效应}（isotope effect）。如果考虑改变每个离子的质量 $M$ 而不改变原子间的作用力，那么我们预期声速按 $M^{-1/2}$ 变化。于是我们预期 $\Theta$ 也按同样的方式变化，从而
\begin{equation*}
T_{c}\propto M^{-1/2},
\tag{11.23}
\end{equation*}
这在许多情形中与观察结果大体相符。

关于 Cooper 对还有一点：如果两个电子除动量相反之外自旋也相反，它们将束缚得更牢。这一点从我们刚才给出的论证看并不明显，而是在计算吸引势 (11.2) 时显露出来的。§11.1 的论证假定我们处理的是可分辨粒子，即自旋相反的电子。倘若两个电子自旋相同，我们就必须援引 Pauli 原理，在计算有效相互作用时使用反对称化的波函数。“交换”项的计入会削弱吸引力，因为这时我们必须设法阻止两个电子进入同一状态或占据同一位置。Cooper 对的单态就是这个准分子的最低态。

\section{超导基态}

上述结果提示了超导态的两个重要特征：它是一个无法用微扰展开中的有限多项来达到的态，因为函数 (11.20) 在 $V=0$ 附近不是 $V$ 的解析函数；它是一个电子以某种方式配成对的态，$\mathbf{k}$ 与 $-\mathbf{k}$ 相配。这个态的构造最初是由 Bardeen、Cooper 和 Schrieffer 所完成的。为了演示他们的解，我们将采用 Bogoliubov 的方法。

完全可以避开有关波函数的一切讨论，而把整个问题用算符语言表述出来，这是非常方便的。我们需要表述这样一个事实：我们有 $N$ 个电子，其能量主要是动能，处在由矢量 $\mathbf{k}$ 和自旋指标 $\sigma$ 标记的态上，并且它们通过势 $V(\mathbf{K})$ 发生相互作用。我们定义湮没算符和产生算符 $b_{\mathbf{k}\sigma}$、$b^{*}_{\mathbf{k}\sigma}$，它们类似于式 (2.129) 和 (10.109) 中定义的算符 $a_{\mathbf{q}}$、$a^{*}_{\mathbf{q}}$，只不过我们现在处理的是费米子。这意味着它们必须服从\emph{反对易关系}（anticommutation relations）
\begin{align*}
[b_{\mathbf{k}\sigma},\,b^{*}_{\mathbf{k}'\sigma'}]_{+} &\equiv b_{\mathbf{k}\sigma}b^{*}_{\mathbf{k}'\sigma'}+b^{*}_{\mathbf{k}'\sigma'}b_{\mathbf{k}\sigma}\\
&=\delta_{\mathbf{k}\mathbf{k}'}\,\delta_{\sigma\sigma'}.
\tag{11.24}
\end{align*}
这个定义的结果是，算符
\begin{equation*}
n_{\mathbf{k}\sigma}=b^{*}_{\mathbf{k}\sigma}b_{\mathbf{k}\sigma}
\tag{11.25}
\end{equation*}
量度态 $(\mathbf{k},\sigma)$ 的占据数，并且如 Pauli 原理所要求的那样，本征值取 0 和 1。当式 (11.24) 中的 $\mathbf{k}$ 与 $\mathbf{k}'$、或 $\sigma$ 与 $\sigma'$ 不相同时，反对易子体现的是反对称原理：交换两个态的标记将使波函数反号。

我们这个系统的 Hamilton 量可以写成
\begin{equation*}
\mathcal{H}=\sum_{\mathbf{k}\sigma}\xi(\mathbf{k})\,b^{*}_{\mathbf{k}\sigma}b_{\mathbf{k}\sigma}+\sum_{\mathbf{k}_1,\,\mathbf{k}_2,\,\mathbf{K}}V(\mathbf{K})\,b^{*}_{\mathbf{k}_2-\mathbf{K},\downarrow}\,b^{*}_{\mathbf{k}_1+\mathbf{K},\uparrow}\,b_{\mathbf{k}_1,\uparrow}\,b_{\mathbf{k}_2,\downarrow}.
\tag{11.26}
\end{equation*}
第一项代表电子的动能，像式 (11.12) 中那样，从费米能级算起。为了方便，常常假定这个能量可以变动——即电子总数 $n$ 并非严格固定——办法是引进电子气的化学势 $\zeta$，并把第一项写成
\begin{align*}
\mathcal{H}_{0}&=\sum_{\mathbf{k}\sigma}n_{\mathbf{k}\sigma}\,\xi(\mathbf{k})\\
&=\sum_{\mathbf{k}\sigma}n_{\mathbf{k}\sigma}\{\mathcal{E}(\mathbf{k})-\zeta\}\\
&=\sum_{\mathbf{k}\sigma}n_{\mathbf{k}\sigma}\,\mathcal{E}(\mathbf{k})-n\zeta.
\tag{11.27}
\end{align*}
式 (11.26) 中的第二项代表由相互作用 (11.2) 所定义的散射，其效果是把 $\mathbf{k}_1$ 和 $\mathbf{k}_2$ 中的电子“湮没”掉，然后再把它们“重新产生”在 $\mathbf{k}_1+\mathbf{K}$ 和 $\mathbf{k}_2-\mathbf{K}$ 处，矩阵元为 $V(\mathbf{K})$。这里我们假定只有自旋相反的电子之间才有相互作用。

Bogoliubov 方法是对算符集 $b_{\mathbf{k}\sigma}$ 与 $b^{*}_{\mathbf{k}\sigma}$ 作一个正则变换，把它们变成具有同样对易关系的一套新的湮没算符和产生算符。特别地，这些新算符必须显示出某种配对性质——正如 Cooper 效应所提示的，它们必须把态 $(\mathbf{k},\uparrow)$ 与态 $(-\mathbf{k},\downarrow)$ 联结起来。

为简单起见，让我们略去自旋指标，写成
\begin{equation*}
\left.\begin{array}{ll}
\beta^{*}_{\mathbf{k}}=u_{\mathbf{k}}b^{*}_{\mathbf{k}}-v_{\mathbf{k}}b_{-\mathbf{k}}, &\quad \beta_{\mathbf{k}}=u_{\mathbf{k}}b_{\mathbf{k}}-v_{\mathbf{k}}b^{*}_{-\mathbf{k}},\\
\beta^{*}_{-\mathbf{k}}=u_{\mathbf{k}}b^{*}_{-\mathbf{k}}+v_{\mathbf{k}}b_{\mathbf{k}}, &\quad \beta_{-\mathbf{k}}=u_{\mathbf{k}}b_{-\mathbf{k}}+v_{\mathbf{k}}b^{*}_{\mathbf{k}}.
\end{array}\right\}
\tag{11.28}
\end{equation*}
（译注：原书 (11.28) 第一式排印为“$\beta_{\mathbf{k}}\ \ u_{\mathbf{k}}b^{*}_{\mathbf{k}}-v_{\mathbf{k}}b_{-\mathbf{k}}$”，漏印等号与 $\beta$ 上的星号；按变换关系此处应为第二式的 Hermite 共轭，译文已补正。）

如果
\begin{equation*}
u^{2}_{\mathbf{k}}+v^{2}_{\mathbf{k}}=1
\tag{11.29}
\end{equation*}
则算符 $\beta_{\mathbf{k}}$、$\beta^{*}_{\mathbf{k}}$ 就会具有费米子应有的反对易关系。
这个变换确实与反铁磁自旋波的变换 (10.130) 本质上相同，只不过现在我们是在 $\mathbf{k}$ 与 $-\mathbf{k}$ 的空间中作一个“真实转动”，于是
\begin{equation*}
u_{\mathbf{k}}=\cos\theta_{\mathbf{k}},\quad v_{\mathbf{k}}=\sin\theta_{\mathbf{k}}
\tag{11.30}
\end{equation*}
满足式 (11.29)。这是因为我们现在处理的是费米子，而不是玻色子。

当我们从式 (11.28) 解出原来的算符并代入 Hamilton 量 (11.26) 时，会得到形形色色的项，其中含有新算符 $\beta_{\mathbf{k}}$、$\beta^{*}_{\mathbf{k}}$ 的乘积。现在，只要某一项中含有湮没算符 $b_{\mathbf{k}}$，我们就可以利用反对易子把它移到右边，例如
\begin{equation*}
\beta_{\mathbf{k}}\beta^{*}_{\mathbf{k}'}=\delta_{\mathbf{k}\mathbf{k}'}-\beta^{*}_{\mathbf{k}'}\beta_{\mathbf{k}}.
\tag{11.31}
\end{equation*}
假定我们的超导态就是这些新算符的“真空”$|0\rangle$。按定义，真空中不存在任何能被湮没的对象，因此
\begin{equation*}
\beta_{\mathbf{k}}|0\rangle=0,
\tag{11.32}
\end{equation*}
对所有 $\mathbf{k}$ 成立。Hamilton 量中的这类项对基态 $|0\rangle$ 的能量没有贡献；于是基态是 $\mathcal{H}$ 中这一部分的本征态。

但是还剩下一些项不能用这种方式消去——即只含有像 $\beta^{*}_{\mathbf{k}}$ 这样的产生算符的项。它们来自 (11.26) 的动能部分，也来自相互作用项的约化——例如当式 (11.31) 中的 $\mathbf{k}$ 与 $\mathbf{k}'$ 恰好相等时。我们发现，Hamilton 量可以写成
\begin{align*}
\mathcal{H}&=2\sum_{\mathbf{k}}\xi(\mathbf{k})\,v^{2}_{\mathbf{k}}+\sum_{\mathbf{k},\,\mathbf{K}}V(\mathbf{K})\,u_{\mathbf{k}}v_{\mathbf{k}}u_{\mathbf{k}+\mathbf{K}}v_{\mathbf{k}+\mathbf{K}}\\
&\quad+\sum_{\mathbf{k}}\bigl\{2\xi(\mathbf{k})\,u_{\mathbf{k}}v_{\mathbf{k}}+(u^{2}_{\mathbf{k}}-v^{2}_{\mathbf{k}})\sum_{\mathbf{K}}V(\mathbf{K})\,u_{\mathbf{k}+\mathbf{K}}v_{\mathbf{k}+\mathbf{K}}\bigr\}\beta^{*}_{\mathbf{k}}\beta^{*}_{-\mathbf{k}}.
\tag{11.33}
\end{align*}
前面的这些项将给出我们基态的能量——只要我们能把最后那个求和中的产生算符消去。

这跟我们在 (10.132) 中遇到的情况非常相似——我们采用同样的手法。我们这样选取变换系数 $u_{\mathbf{k}}$、$v_{\mathbf{k}}$，使每个“危险”项都等于零。令
\begin{equation*}
2\xi(\mathbf{k})\,u_{\mathbf{k}}v_{\mathbf{k}}+(u^{2}_{\mathbf{k}}-v^{2}_{\mathbf{k}})\sum_{\mathbf{K}}V(\mathbf{K})\,u_{\mathbf{k}+\mathbf{K}}v_{\mathbf{k}+\mathbf{K}}=0.
\tag{11.34}
\end{equation*}
一般说来，这是关于未知函数 $u_{\mathbf{k}}$ 和 $v_{\mathbf{k}}$ 的一个复杂的积分方程，而且它们还须满足条件 (11.29)。但让我们采用 (11.8) 和 (11.9) 的假定，把吸引势 $V(\mathbf{K})$ 简化为在能量范围 $\pm w$ 之内取常数值 $V$、在此范围之外为零。我们把式 (11.34) 中对 $\mathbf{K}$ 的求和记作一个参量
\begin{equation*}
\Delta_{0}=-V\sum_{-w}^{w}u_{\mathbf{k}+\mathbf{K}}v_{\mathbf{k}+\mathbf{K}},
\tag{11.35}
\end{equation*}
然后联立求解 (11.34) 与 (11.29)。结果是
\begin{equation*}
u^{2}_{\mathbf{k}}=\frac{1}{2}\biggl[1+\frac{\xi(\mathbf{k})}{\sqrt{\{\Delta^{2}_{0}+\xi^{2}(\mathbf{k})\}}}\biggr],\quad v^{2}_{\mathbf{k}}=\frac{1}{2}\biggl[1-\frac{\xi(\mathbf{k})}{\sqrt{\{\Delta^{2}_{0}+\xi^{2}(\mathbf{k})\}}}\biggr],
\tag{11.36}
\end{equation*}
把它代回 (11.35)，便得到
\begin{align*}
1&=-\frac{1}{2}V\sum_{-w}^{w}\{\Delta^{2}_{0}+\xi^{2}(\mathbf{k})\}^{-\frac{1}{2}}\\
&=-\frac{1}{2}\,\mathcal{N}(\mathcal{E}_{F})\,V\ln\frac{2w}{\Delta_{0}},
\tag{11.37}
\end{align*}
这是在费米面附近对 $\xi$ 的范围积分而来的（记住自旋只计一种）。这与 (11.19) 完全相同。参量 $\Delta_{0}$ 恰恰就是一个 Cooper 对的结合能，与 (11.20) 或 (11.22) 中的一样。

我们可以把这个结果代回 (11.33) 的其余部分——现在其中所有的算符都已消去。我们得到
\begin{equation*}
\mathcal{E}_{0}=2\sum_{\xi<0}\xi(\mathbf{k})-\frac{1}{2}\Delta^{2}_{0}\sum_{\mathbf{k}}\{\Delta^{2}_{0}+\xi^{2}(\mathbf{k})\}^{-\frac{1}{2}}.
\tag{11.38}
\end{equation*}
于是，态 $|0\rangle$ 是 $\mathcal{H}$ 的一个本征态——但其能量低于电子填满的费米球，后者的能量本来正好是各被占据态的 $\xi(\mathbf{k})$ 值之和，如 (11.27) 那样。还可以证明，对于电子总数的期望值 $n$ 的变动而言，这个态使能量取极小，正如 (11.27) 所要求的。

\section{准粒子与能隙}

超导态 $|0\rangle$ 的本性是什么？按定义，它是算符 $\beta^{*}_{\mathbf{k}}$、$\beta_{\mathbf{k}}$ 的“真空”。这些算符能够产生和消灭真空的激发，正如原来的算符 $b^{*}_{\mathbf{k}}$ 和 $b_{\mathbf{k}}$ 能够“产生”和“消灭”电子一样。这样一种激发看上去会是什么样子？按照 (11.28)，$\beta^{*}_{\mathbf{k}}$ 的效果是以振幅 $u_{\mathbf{k}}$“产生”一个处在态 $\mathbf{k}$ 中的电子，同时又以振幅 $v_{\mathbf{k}}$“消灭”一个处在态 $-\mathbf{k}$ 中的电子。

这些系数由 (11.36) 给出。假定 $\mathbf{k}$ 远在费米面之上，使得 $\xi(\mathbf{k})\gg\Delta_{0}$，那么 $u^{2}_{\mathbf{k}}\approx 1$，而 $v^{2}_{\mathbf{k}}\approx 0$。我们的激发看上去就会像一个普通的电子。
\begin{figure}[H]
\centering
\includegraphics[width=0.72\textwidth]{fig_203.pdf}
\caption*{图 203\quad 准粒子。}
\end{figure}
但是现在，当我们逼近费米能级时，$\xi(\mathbf{k})$ 趋于零，$b^{*}_{\mathbf{k}}$ 的振幅将减小，而 $b_{-\mathbf{k}}$ 的振幅将增大。当 $\xi(\mathbf{k})$ 为负时，$v^{2}_{\mathbf{k}}$ 的值将大于 $\frac{1}{2}$；当我们远在费米能级之下时，$v^{2}_{\mathbf{k}}\approx 1$ 而 $u^{2}_{\mathbf{k}}\approx 0$。这时我们的激发相应于在 $-\mathbf{k}$ 处一个电子的“毁灭”，也就是说，它是本已被占据的诸态中的一个“空穴”。

因此，超导态的激发是颇为奇特的\emph{准粒子}（quasi-particle），当它们通过费米能级时，会从“电子”变成“空穴”。在 $\Delta_{0}$ 附近的能量范围内，每个准粒子都是 $\mathbf{k}$ 处的一个电子与 $-\mathbf{k}$ 处的一个空穴的混合体。Cooper 对本身不再出现，但是动量相反（当然还有自旋相反）的电子之间的关联是显而易见的。实际上，我们是在坚持：配对态中各波函数的相位之间应当存在一种关系。

激发的能量也可以计算。我们在 (11.26) 的展开式中寻找含下列算符
\begin{equation*}
n'_{\mathbf{k}}=\beta^{*}_{\mathbf{k}}\beta_{\mathbf{k}}
\tag{11.39}
\end{equation*}
的各项的系数，因为它量度的是在此模式中被激发的准粒子的数目。结果是
\begin{align*}
\epsilon(\mathbf{k})&=\xi(\mathbf{k})\,(u^{2}_{\mathbf{k}}-v^{2}_{\mathbf{k}})-2u_{\mathbf{k}}v_{\mathbf{k}}\sum_{\mathbf{K}}V(\mathbf{K})\,u_{\mathbf{k}+\mathbf{K}}v_{\mathbf{k}+\mathbf{K}}\\
&=\sqrt{\{\Delta^{2}_{0}+\xi^{2}(\mathbf{k})\}}
\tag{11.40}
\end{align*}
（对我们的简化相互作用，用了 (11.35) 和 (11.36)）。
\begin{figure}[H]
\centering
\includegraphics[width=0.72\textwidth]{fig_204.pdf}
\caption*{图 204\quad 准粒子的能量。}
\end{figure}
这一点非常重要，因为它证明了在超导基态之上存在\emph{能隙}（energy gap）$\Delta_{0}$。当 $\xi(\mathbf{k})\gg\Delta_{0}$ 时，$\epsilon(\mathbf{k})\approx\xi(\mathbf{k})$——准粒子的能量就是把一个普通电子激发到费米能级之上所需的能量。但当 $\mathbf{k}$ 趋近费米面（那里 $\xi(\mathbf{k})$ 为零）时，准粒子的能量趋于常数值 $\Delta_{0}$。然后，随着 $\mathbf{k}$ 再进一步减小，$\epsilon(\mathbf{k})$ 又重新增大，直至几乎等于 $|\xi(\mathbf{k})|$。这当然就是把一个电子从负能量态 $\xi(\mathbf{k})$ 中移出所需的能量，正如我们根据对准粒子算符的解释所预期的那样。
因此，超导态是一种“凝聚”态，其含义是：要使整个系统产生一个激发态，需要有限的能量 $\Delta_{0}$。在某些方面，它有点像半导体——在半导体中，要想把一个电子送入导带，也必须越过一有限的带隙。这个类比有助于粗略地解释超导体的某些特殊性质，例如电磁吸收（参见 §8.5）和隧穿（§6.8）。

\section{能隙随温度的变化}

基态只适用于 $T=0$ 的情形。在有限温度下，我们预期会有准粒子按通常的 费米–狄拉克 函数被激发：
\begin{equation*}
f_{\mathbf{k}}=\frac{1}{e^{\epsilon(\mathbf{k})/kT}+1}.
\tag{11.41}
\end{equation*}
在正常态中，这类激发无论在费米能级之上还是之下，都彼此独立。在超导态中，它们倾向于相互合作地发生作用，并破坏能隙。

让我们假定所处理的是态 $|f_{\mathbf{k}}\rangle$，其中第 $\mathbf{k}$ 个态中准粒子的平均数目由 (11.41) 给出。湮没算符 $b_{\mathbf{k}}$ 作用在这个态上的效果不再是零。平均而言，我们可以把 (11.25) 换成
\begin{equation*}
b^{*}_{\mathbf{k}}b_{\mathbf{k}}|f_{\mathbf{k}}\rangle=f_{\mathbf{k}}|f_{\mathbf{k}}\rangle.
\tag{11.42}
\end{equation*}
把 Hamilton 量约化成对角形式的做法，当应用于态 $|0\rangle$ 时，依靠的是 (11.32)。如果现在对态 $|f_{\mathbf{k}}\rangle$ 试着做同样的事情，就会多出一些额外的项；它们来自像 (11.31) 那样把算符乘积约化成标准次序的过程，并乘在 (11.33) 中那些“危险”乘积 $\beta^{*}_{\mathbf{k}}\beta^{*}_{-\mathbf{k}}$ 上。为了从 Hamilton 量中去掉这些“产生电子对”的部分，我们的条件 (11.34) 必须改成为
\begin{equation*}
2\xi(\mathbf{k})\,u_{\mathbf{k}}v_{\mathbf{k}}+(u^{2}_{\mathbf{k}}-v^{2}_{\mathbf{k}})\sum_{\mathbf{K}}V(\mathbf{K})\,u_{\mathbf{k}+\mathbf{K}}v_{\mathbf{k}+\mathbf{K}}(1-2f_{\mathbf{k}+\mathbf{K}})=0
\tag{11.43}
\end{equation*}
（这里计及与每个 $\mathbf{k}$ 相联系的两个可能的自旋态）。

对于简化的相互作用 $V$，这个积分方程可以这样求解：把 (11.35) 和 (11.36) 中的 $\Delta_{0}$ 换成
\begin{equation*}
\Delta=-V\sum_{-w}^{w}u_{\mathbf{k}+\mathbf{K}}v_{\mathbf{k}+\mathbf{K}}(1-2f_{\mathbf{k}+\mathbf{K}}).
\tag{11.44}
\end{equation*}
代替 (11.37)，利用 (11.41)，我们便得到
\begin{align*}
1&=-\frac{1}{2}V\sum_{-w}^{w}\frac{1}{\epsilon(\mathbf{k})}\tanh\frac{\epsilon(\mathbf{k})}{2kT}\\
&=-\frac{1}{2}\,\mathcal{N}(\mathcal{E}_{F})\,V\int_{-w}^{w}\frac{1}{\epsilon}\tanh\frac{\epsilon}{2kT}\,d\xi,
\tag{11.45}
\end{align*}
其中 $\epsilon$ 和以前一样，是一个激发的能量，
\begin{equation*}
\epsilon(\mathbf{k})=\sqrt{\{\Delta^{2}+\xi^{2}(\mathbf{k})\}}.
\tag{11.46}
\end{equation*}
式 (11.45) 与式 (11.46) 这两个方程给出 $\Delta$ 与 $T$ 之间的一个隐含关系。这个关系相当杂乱，但要点在于：当 $T$ 增大时，$\Delta$ 从 $T=0$ 处的 $\Delta_{0}$ 逐渐减小。换言之，能隙 $\Delta$ 随温度升高而减小，并在一个完全确定的温度 $T_{c}$ 处闭合为零；由 (11.46)，在该处 $\epsilon(\mathbf{k})=|\xi(\mathbf{k})|$。于是，我们可以通过积分
\begin{equation*}
\frac{2}{\mathcal{N}(\mathcal{E}_{F})\,V}=2\int_{0}^{w}\frac{1}{\epsilon}\tanh\left(\frac{\epsilon}{2kT_{c}}\right)d\epsilon,
\tag{11.47}
\end{equation*}
求出 $T_{c}$，其解为
\begin{equation*}
\Delta_{0}\approx 1.76\,kT_{c},
\tag{11.48}
\end{equation*}
\begin{figure}[H]
\centering
\includegraphics[width=0.72\textwidth]{fig_205.pdf}
\caption*{图 205\quad 能隙随温度的变化。}
\end{figure}
式中的 $\Delta_{0}$ 由 (11.22) 表示，即 $T=0$ 时的能隙。这样一来，§11.2 中由 Cooper 对的结合能来估计 $T_{c}$ 的讨论就得到了印证。

在 $T_{c}$ 以上，方程没有解。恰好低于 $T_{c}$ 时，能隙从零陡峭地升起，其变化规律为
\begin{equation*}
\Delta(T)\approx 3.2\,kT_{c}\left(1-\frac{T}{T_{c}}\right)^{\frac{1}{2}},
\tag{11.49}
\end{equation*}
这可以通过在 (11.47) 之解的邻域内寻求 (11.45) 与 (11.46) 的近似解来验证。在更低的温度下，它增大得较为缓慢，在大约 $T\sim\frac{1}{2}T_{c}$ 以下逐渐变平而趋于 $\Delta_{0}$。

从对我们的准粒子系统的能量所作的仔细计算——其中计及能隙随温度
%（句子未完，下接书页 394 / p408 页首 "ture, one can write down formulae for the specific heat."——应由下一部分的 B-TAIL 以 %JOIN% 续接"的变化，便可以写下比热的公式。……"；式 (11.50)、(11.51)、图 206 及 §11.6 均在 p408 起，归下一部分。）
 度的变化，即可写出比热的公式。恰在转变温度处，比热有一个数值为
\begin{equation*}
C_{es} = 1.43\,\gamma T_c, \tag{11.50}
\end{equation*}
的跃变，其中 $\gamma$ 是正常态电子比热公式 (4.46) 中的系数。在更低的温度下，能隙趋于在比热中占据主导，但像
\begin{equation*}
C_{es} \approx e^{-\Delta(T)/T} \tag{11.51}
\end{equation*}
这样简单的公式并不适用，除非 $\Delta$ 已经稳定到其极限值 $\Delta_0$。

\begin{figure}[H]
\centering
\includegraphics[width=0.72\textwidth]{fig_206.pdf}
\caption*{图 206\quad 超导体的比热，与正常金属电子比热的比较。}
\end{figure}

\section{持续电流}

超导态最引人注目的特征是，金属对于电流的流动似乎完全不表现出任何电阻。要在 BCS 理论中理解这一点，我们注意到能隙并不系结于晶格——它不像半导体中那样固定在倒格子空间的区界上——而只是出现在费米能级处。我们上面所用的论证并不依赖于费米面是球形的；对于任意形状的费米面，它也应当成立。

如 §7.2 所示，电场的作用是使费米分布相对于倒格子整体位移一个恒定矢量。电流 $\mathbf{J}$ 对应于自由电子球的如下位移：
\begin{equation*}
\boldsymbol{\delta}\mathbf{k} = \frac{m}{ne\hbar}\,\mathbf{J} \tag{11.52}
\end{equation*}
能隙并不妨碍这一移动；它随费米面一起被带到新的位置。

要用 BCS 理论的语言描述载流状态，我们只需假设每个电子对的质心移动了 $\boldsymbol{\delta}\mathbf{k}$。在构造式 (11.28) 的准粒子算符 $\beta_{\mathbf{k}}$、$\beta_{\mathbf{k}}^{*}$ 时，我们如今把态 $(\mathbf{k}+\boldsymbol{\delta}\mathbf{k},\uparrow)$ 与态 $(-\mathbf{k}+\boldsymbol{\delta}\mathbf{k},\downarrow)$ 相配；论证的其余部分照旧不变。

\begin{figure}[H]
\centering
\includegraphics[width=0.72\textwidth]{fig_207.pdf}
\caption*{图 207\quad (a)半导体中，能隙系结于布里渊区。(b)超导体中，能隙由费米面携带。}
\end{figure}

唯一需要指出的一点是，这样一种电子分布的总动能将增加
\begin{equation*}
\delta\mathcal{E} \approx \frac{\hbar^{2}(\delta k)^{2}}{2m} \tag{11.53}
\end{equation*}
的量（每个电子）。如果这个能量超过了超导态的结合能 (11.38)，我们就不能指望系统仍凝聚在能隙之下。这就为所能承载的电流设定了一个极限。

但载流态最重要的性质是：一旦建立起来，它就不容易被破坏。如 §7.5 所讨论的，阻碍正常电流的那些过程是“单电子”跃迁——单个粒子通过与声子或杂质的相互作用，从费米面的一侧被散射到另一侧。在超导态中，这样的过程只能通过产生准粒子来进行——在电子原来所在处（比如说 $\mathbf{k}+\mathbf{q}$）产生一个空穴，而电子又在 $(-\mathbf{k}+\mathbf{q})$ 处重新出现。这需要能量 $\Delta$，而单个声子提供不了这份能量。要把电流降到零，就得设法让所有电子同时慢下来，使整个位移了的费米面平滑地退回原点。对这样一种可能性极小的过程，并不存在简单的机制，因此在一段环形的超导材料中，环行电流一旦建立，几乎可以永不衰减地持续下去。

BCS 理论似乎关键性地依赖于存在明晰定义的 Bloch 态 $|\mathbf{k},\uparrow\rangle$ 和 $|-\mathbf{k},\downarrow\rangle$。然而，超导电性却常常在极无序、“肮脏”的合金中被观察到；在这些合金里，电子作普通导电时的平均自由程短得使我们无法把 $\mathbf{k}$ 当作好的量子数来处理。真实的情况是：超导“配对”可以在任何两个经由时间反演相互关联的态之间发生（§3.11）。全部必要条件不过是：电子在反转其自旋之后，能够沿原路折返它穿过材料的路径，无论这条路径多么迂回曲折。这样一来，Fröhlich 相互作用的理论看上去会非常复杂，但只要系统的本征态能够按时间反演配对来分类，且费米能处于态密度高的区域，那么在原则上没有任何东西妨碍一个宏观超导凝聚相的出现。基于同样的理由，把近独立电子态重新诠释为多体准粒子激发（§5.8），也不构成 BCS 理论的障碍。

然而，如果材料中含有内建的磁场，例如在铁磁体中，时间反演对称性的假设便不再成立（§9.3），因为我们无法假设在让单个电子沿其路径折返的假想过程中，这些磁场也会一并被反转。因此，相当低的磁性杂质浓度（§10.6）就足以摧毁超导电性。在杂质浓度稍低一些时，则可以观察到一种称为无能隙超导电性（gapless superconductivity）的状态——表面上是超导的、宏观上是有序的（§11.9），但准粒子谱中却没有实际的能隙 (11.40)。

\section{London 方程}

超导体在外磁场中的行为颇为奇特。其中一部分行为表面上可以理解为电导率无限的后果。我们无法在超导体内部建立起磁场，因为在磁场建立时，样品表面会感应出持续电流（persistent current），从而把内部屏蔽起来。但这个解释并不足以说明全部现象。

我们需要考察由具有矢势 $\mathbf{A}(\mathbf{r})$ 的磁场感应出的电流。按照式 (10.3)，这会引起动能算符的改变；在态 $\psi(\mathbf{r})$ 中，这个改变的期望值为
\begin{equation*}
\psi^{*}\left\{\frac{1}{2m}\left(\frac{\hbar}{i}\nabla - \frac{e}{c}\mathbf{A}\right)^{2} - \left(-\frac{\hbar^{2}}{2m}\nabla^{2}\right)\right\}\psi = -\frac{1}{c}\,\mathbf{J}(\mathbf{r})\cdot\mathbf{A}(\mathbf{r}), \tag{11.54}
\end{equation*}
在这里，我们是把这个局部的“动能密度”等同于密度为 $\mathbf{J}(\mathbf{r})$ 的电流在矢势 $\mathbf{A}(\mathbf{r})$ 中运动时相应的经典能量密度。于是，电流密度就是 (11.54) 对 $\mathbf{A}(\mathbf{r})$ 的导数，其形式为
\begin{equation*}
\mathbf{J}(\mathbf{r}) = \frac{e\hbar}{2m}\,\frac{1}{i}\left(\psi^{*}\nabla\psi - \psi\nabla\psi^{*}\right) - \frac{e^{2}}{mc}\,\psi^{*}\psi\,\mathbf{A}. \tag{11.55}
\end{equation*}
第一项就是大家熟知的无磁场时的电流公式。

在普通金属中，静场下这两部分贡献几乎相消，只剩下微弱的抗磁性（参见 §10.1）。然而在超导体中，“波函数” $\psi$ 的那些态（将在 §11.9 中更详细地讨论）具有某种“刚性”。举例来说，假设磁场的施加对波函数毫无影响，$\psi$ 保持等于零场中的值 $\psi_0$，那么我们可以肯定
\begin{align*}
\mathbf{J}_0(\mathbf{r}) &= \frac{e\hbar}{2mi}\left(\psi_0^{*}\nabla\psi_0 - \psi_0\nabla\psi_0^{*}\right)\\
&= 0. \tag{11.56}
\end{align*}
于是在超导体中，剩下的便是
\begin{align*}
\mathbf{J}(\mathbf{r}) &= -\frac{e^{2}}{mc}\,\psi^{*}(\mathbf{r})\psi(\mathbf{r})\,\mathbf{A}(\mathbf{r})\\
&= -\frac{e^{2}n}{mc}\,\mathbf{A}(\mathbf{r}), \tag{11.57}
\end{align*}
其中 $n$ 是当地的电子密度。

这个不再出现假想“波函数”的公式，就是著名的 London 方程（London equation）；它对解释超导体的许多性质大有帮助。例如，我们知道电导率是无限的，因此样品内部没有电场。静磁场将通过 Maxwell 方程
\begin{equation*}
\nabla\wedge\mathbf{H} = \frac{4\pi}{c}\,\mathbf{J}. \tag{11.58}
\end{equation*}
与电流 $\mathbf{J}$ 相联系。但 (11.57) 可以写成
\begin{align*}
\mathbf{H} &= \nabla\wedge\mathbf{A}\\
&= -\frac{mc}{e^{2}n}\,\nabla\wedge\mathbf{J}. \tag{11.59}
\end{align*}
把 (11.58) 与 (11.59) 结合起来，我们得到
\begin{equation*}
\nabla^{2}\mathbf{H} = \frac{1}{\lambda^{2}}\,\mathbf{H}, \qquad \nabla^{2}\mathbf{J} = \frac{1}{\lambda^{2}}\,\mathbf{J}, \tag{11.60}
\end{equation*}
其中
\begin{equation*}
\lambda = \sqrt{\frac{mc^{2}}{4\pi ne^{2}}}. \tag{11.61}
\end{equation*}
对于恰在超导体表面之外并平行于表面的磁场 $\mathbf{H}_0$，(11.60) 的解是
\begin{equation*}
\mathbf{H} = \mathbf{H}_0\,e^{-z/\lambda}, \tag{11.62}
\end{equation*}
其中 $z$ 是深入样品内部量度的距离。这样，磁场迅速衰减，只能穿透 $\lambda$ 的距离——这就是所谓的穿透深度（penetration depth），其量级为 $5\times10^{-6}$ cm。

这一结果并不依赖于磁场是如何建立的。对于在磁场中从转变温度之上冷却下来的样品，它同样成立。当金属变为超导体时，磁场被排了出去。这就是 Meissner 效应（Meissner effect），它无法仅凭无限电导率来解释。

超导态要求时间反演对称性，而且是“完全抗磁的”，这一事实提示：足够强的磁场应当能够改变这个态。确实，当施加的磁场超过临界值 $H_c$ 时，块体金属中的超导电性即被破坏。对第一类超导体（type I superconductor，见 §11.10），这个场可以由比热曲线通过热力学论证计算出来。临界场处的场能 $H_c^{2}/8\pi$ 必须等于超导态与正常态之间的自由能之差。举例来说，假设超导态的比热正比于 $T^{3}$，而正常态具有普通的电子比热 $\gamma T$，那么可以证明临界场应是温度的函数：
\begin{equation*}
H_c = \sqrt{(2\pi)\gamma T_c\{1-(T/T_c)^{2}\}}, \tag{11.63}
\end{equation*}
其中 $T_c$ 是临界温度 (11.48)。实际上，正如在 (11.51) 处所指出的，超导体的比热是温度的复杂函数，所以这个公式只是近似正确。

值得注意的是，这一论证假设我们能够把\emph{全部}磁场排除在样品之外。对于厚度与 $\lambda$ 相当的很小或很薄的金属块，磁场会穿透进去，热力学论证便不再有效。这时超导态可以存在于比 (11.63) 算出的值高得多的磁场之下。第二类超导体（type II superconductor）中发生的正是这种情形：在那里，即使在块体样品中，这类细丝也可以通过与位错、杂质等的相互作用而稳定下来（见 §11.10）。

\section{相干长度}

推导 London 方程 (11.57) 时的假设是：超导波函数是“刚性的”，不受磁场施加的影响。这显然不正确；我们知道，足够强的磁场会破坏超导态。我们需要计算 (11.55) 中的波函数 $\psi$ 在微扰 $\mathbf{A}(\mathbf{r})$ 作用于超导基态时所产生的改变。

这在 §11.3 的 Bogoliubov 形式框架内并不难算出。假设我们考虑矢势的一个单独的 Fourier 分量
\begin{equation*}
\mathbf{A}(\mathbf{r}) = \mathbf{A}_{\mathbf{q}}\,e^{i\mathbf{q}\cdot\mathbf{r}}. \tag{11.64}
\end{equation*}
在 $\exp(i\mathbf{k}\cdot\mathbf{r})$ 和 $\exp(i\mathbf{k}'\cdot\mathbf{r})$ 这两个普通自由电子态之间，这一项之一的矩阵元将是
\begin{equation*}
\int e^{-i\mathbf{k}'\cdot\mathbf{r}}\,\mathbf{A}(\mathbf{r})\,e^{i\mathbf{k}\cdot\mathbf{r}}\,\mathrm{d}\mathbf{r} = \mathbf{A}_{\mathbf{q}}\,\delta(\mathbf{k}+\mathbf{q}-\mathbf{k}'). \tag{11.65}
\end{equation*}
这样，用湮没和产生算符 (11.24) 的语言来说，$\mathbf{A}(\mathbf{r})$ 的行为就如同算符
\begin{equation*}
\mathbf{A} = \mathbf{A}_{\mathbf{q}}\,b_{\mathbf{k}+\mathbf{q}}^{*}b_{\mathbf{k}} \tag{11.66}
\end{equation*}
一样。精确到 $\mathbf{A}$ 的一阶，磁场在 Hamiltonian 中产生的实际微扰就是算符
\begin{equation*}
-\frac{e\hbar}{2mci}\left(\nabla\cdot\mathbf{A}+\mathbf{A}\cdot\nabla\right) = -\frac{e\hbar}{2mc}\sum_{\mathbf{k}}\left\{(\mathbf{k}+\mathbf{q})+\mathbf{k}\right\}\cdot\mathbf{A}_{\mathbf{q}}\,b_{\mathbf{k}+\mathbf{q}}^{*}b_{\mathbf{k}}. \tag{11.67}
\end{equation*}
这很容易通过与 (11.65) 和 (11.66) 的类比得到；在所有“裸电子”态之间，右边的算符与左边的算符具有相同的矩阵元。

现在对这个表达式施加变换 (11.28)。用准粒子算符 $\beta_{\mathbf{k}}$、$\beta_{\mathbf{k}}^{*}$ 表出的结果有点复杂，但由于我们是要把微扰作用到超导基态 $|0\rangle$ 上，可以利用 (11.32) 消去所有含有湮没算符的项。剩下的就是含有两个产生算符之积的项；我们的微扰变为
\begin{equation*}
-\frac{e\hbar}{2mc}\sum_{\mathbf{k}}\,(2\mathbf{k}+\mathbf{q})\cdot\mathbf{A}_{\mathbf{q}}\,(v_{\mathbf{k}+\mathbf{q}}u_{\mathbf{k}} - v_{\mathbf{k}}u_{\mathbf{k}+\mathbf{q}})\,\beta_{\mathbf{k}+\mathbf{q}}^{*}\beta_{-\mathbf{k}}^{*}, \tag{11.68}
\end{equation*}
其中 $v_{\mathbf{k}}$、$u_{\mathbf{k}}$ 等是 (11.28) 中的系数。磁场产生净动量为 $\mathbf{q}$ 的准粒子对，正如我们从关于持续电流性质的论证（§11.6）中所应预期的那样。

我们用 (11.68) 作为微扰来计算基态波函数的改变，然后用这个“非刚性”的态代替 $\psi$，重新由 (11.55) 计算电流。结果如下：
\begin{equation*}
\mathbf{J}_{\mathbf{q}} = -\frac{e^{2}n}{mc}\,\mathbf{A}_{\mathbf{q}} + \frac{e^{2}}{2m^{2}c}\sum_{\mathbf{k}}\frac{(2\mathbf{k}+\mathbf{q})\left\{(2\mathbf{k}+\mathbf{q})\cdot\mathbf{A}_{\mathbf{q}}\right\}(v_{\mathbf{k}}u_{\mathbf{k}+\mathbf{q}} - v_{\mathbf{k}+\mathbf{q}}u_{\mathbf{k}})}{\epsilon(\mathbf{k})+\epsilon(\mathbf{k}+\mathbf{q})}, \tag{11.69}
\end{equation*}
其中准粒子对的能量 $\epsilon(\mathbf{k})+\epsilon(\mathbf{k}+\mathbf{q})$ 由 (11.40) 算出。

这一计算中存在一些形式上的困难——关于矢势规范选取的困难，这类困难在磁场中粒子的量子理论中经常出现。这里我们可以假定 $\mathbf{A}(\mathbf{r})$ 是纯横向的，即 $\mathbf{q}\cdot\mathbf{A}_{\mathbf{q}}=0$，从而避开这些困难。为方便起见，把 $v_{\mathbf{k}}$ 等的乘积对称化；我们得到
\begin{equation*}
\mathbf{J}_{\mathbf{q}} = -\frac{e^{2}n}{mc}\,\mathbf{A}_{\mathbf{q}} + \frac{2e^{2}}{m^{2}c}\sum_{\mathbf{k}}\frac{\mathbf{k}(\mathbf{k}\cdot\mathbf{A}_{\mathbf{q}})\left(v_{\mathbf{k}-\frac{1}{2}\mathbf{q}}u_{\mathbf{k}+\frac{1}{2}\mathbf{q}} - v_{\mathbf{k}+\frac{1}{2}\mathbf{q}}u_{\mathbf{k}-\frac{1}{2}\mathbf{q}}\right)}{\epsilon(\mathbf{k}-\frac{1}{2}\mathbf{q})+\epsilon(\mathbf{k}+\frac{1}{2}\mathbf{q})}. \tag{11.70}
\end{equation*}
利用 (11.36) 中 $u_{\mathbf{k}}$ 和 $v_{\mathbf{k}}$ 用能量 $\xi(\mathbf{k})$ 表出的公式，这个和式很容易算出。我们不打算详细追述这一论证；它不过是初等的几何。我们只是把 (11.70) 写成
\begin{equation*}
\mathbf{J}_{\mathbf{q}} = -\frac{c}{4\pi}\,\Gamma(\mathbf{q})\,\mathbf{A}_{\mathbf{q}} \tag{11.71}
\end{equation*}
的形式，用以表示电流与矢势之间的关系；$\Gamma(\mathbf{q})$ 是所施加磁场的波数 $q$ 的一个完全确定的函数。

从 (11.70) 的形式可以相当清楚地看出，当 $q\to 0$ 时和式消失。因此，对缓慢变化的磁场，London 公式 (11.57) 成立；一般说来，Meissner 效应随之而来。但随着 $q$ 增大，(11.70) 中的和式——它本质上是正的——趋于抵消首项。的确，对大的 $q$ 值，我们考虑的是能量远大于能隙的激发，准粒子的行为就像普通的电子和空穴（依其在费米能级之上还是之下而定）。这时可以证明——这是运用该理论的一个很好的练习——London 项恰好被抵消。换言之，当 $q$ 变大时，
\begin{equation*}
\Gamma(q) \to 0 \tag{11.72}
\end{equation*}
$\Gamma(q)$ 在这两个极限之间的实际形式如图 208 所示。很自然可以设想（而且这同样可以由 (11.70) 验证），当 $q$ 大到足以产生一个激发时——也就是当 $q$ 超过
\begin{equation*}
\hbar v_F q_c \sim \Delta_0. \tag{11.73}
\end{equation*}
所给出的 $q_c$ 时——我们将从“London”区制过渡到 (11.72) 所示的区制。

\begin{figure}[H]
\centering
\includegraphics[width=0.72\textwidth]{fig_208.pdf}
\caption*{图 208\quad 响应函数：(a)在波数空间中，(b)在真实空间中。}
\end{figure}

这个公式所定义的量 $1/q_c$，就是在纯金属的理想情形下\emph{相干长度}（coherence length）$\xi$ 的值。

为了看出它的意义，让我们把 (11.71) 表成\emph{真实空间}中电流与矢势之间的关系。它变为
\begin{equation*}
\mathbf{J}(\mathbf{r}) = \int \Gamma(\mathbf{r}-\mathbf{r}')\,\mathbf{A}(\mathbf{r}')\,\mathrm{d}\mathbf{r}', \tag{11.74}
\end{equation*}
其中 $\Gamma(\mathbf{r})$ 是 $\Gamma(\mathbf{q})$ 的 Fourier 变换。在 London 理论中 $\Gamma(\mathbf{q})$ 是常数，所以 $\Gamma(\mathbf{r})$ 是 $\delta$ 函数；$\mathbf{J}(\mathbf{r})$ 的值只依赖于 $\mathbf{A}(\mathbf{r})$ 的\emph{局域}值，如 (11.57) 那样。但 (11.71) 的变换要更复杂。$\Gamma(\mathbf{r})$ 的范围将是相干长度 $\xi$，于是我们将得到 $\mathbf{J}$ 与 $\mathbf{A}$ 之间的一个\emph{非局域}关系。它与反常趋肤效应理论中使用的 Chambers 公式 (8.101) 颇为相似。

这一由 Pippard 从唯象上建立的结果，对于反常趋肤效应、磁场穿透等性质的细致分析十分重要。实际上，$\xi$ 就是一个 Cooper 对的“尺度”；它是配对的电子的相位必须发生关联的距离。在“肮脏”的金属中，当电子平均自由程 $\Lambda$ 小于 $1/q_c$ 时，相干长度 $\xi$ 的表观值不会超过 $\Lambda$ 本身。这时相位关联在一定程度上被破坏，但配对本身并未被摧毁，材料仍是超导体。

有意思的是看一看这种相干性在数学理论中是如何表达的。从 (11.67) 到 (11.68) 的步骤把它显示得非常清楚。算符 (11.67) 引起单电子态之间的散射。在正常金属中，我们应当把这些事件当作相互独立的事件来处理，并把各个矩阵元的平方相加。这样，在 (11.69) 中我们将恰好有对应于
\begin{equation*}
(v_{\mathbf{k}}u_{\mathbf{k}+\mathbf{q}})^{2} + (u_{\mathbf{k}}v_{\mathbf{k}+\mathbf{q}})^{2} \tag{11.75}
\end{equation*}
的项出现在分子中。

但在超导态中，该算符以\emph{单一}的矩阵元
\begin{equation*}
(v_{\mathbf{k}}u_{\mathbf{k}+\mathbf{q}} - u_{\mathbf{k}}v_{\mathbf{k}+\mathbf{q}}) \tag{11.76}
\end{equation*}
产生两个动量近乎相反的激发，再取其平方。这给出的结果与 (11.75) 不同；正是乘积项在微扰所产生的两个准粒子之间引起关联，即相干性。但这一项只对小于 $q_c$ 的 $q$ 才重要；这正是相干长度的意义所在。

\section{非对角长程序}

超流电流可以沿着长达许多千米的回路流动。这样的回路上某一点的物理条件会受到相距极远的其他各点条件的影响。因此，电子系统表现出一种特殊形式的\emph{长程序}（long range order），我们想从数学上刻画它，同时适当地顾及当我们从一个区域走到另一区域时材料、温度或磁场的改变。

在半导体的宏观理论中，这些局域特征被概括为各类载流子的密度和迁移率。类似地，在 Gorter--Casimir 二流体模型（two-fluid model）中，我们假设超流电流由数密度为 $n_s$ 的“超流电子”携带，$n_s$ 依赖于温度等。如果每个这样的电子以速度 $\mathbf{v}_s$ 运动、携带电荷 $e^{*}$（不一定是普通的电子电荷 $e$），那么总的超流电流密度必须是
\begin{equation*}
\mathbf{J}_s = n_s e^{*}\mathbf{v}_s. \tag{11.77}
\end{equation*}
这个模型在 BCS 理论的框架内几乎可以得到论证：把“正常”电子的平衡浓度
\begin{equation*}
n_n = n - n_s, \tag{11.78}
\end{equation*}
认同为 (11.41) 所给出的准粒子激发的浓度。在 $T=0$ 时没有准粒子，故 $n_s=n$；在临界温度 $T_c$ 处，$n_s$ 迅速降到零。我们甚至可以在热力学上逼近图 206 的比热反常，办法是在超流相与“正常”相之间指定一个自由能差。于是我们可以设想，例如，超导隧穿（§11.11）这类空间性的效应，在数学上可以用 $n_s$ 的局域值随位置的变化来描述。

然而这个描述并不充分。正如我们在 §11.8 中看到的，超导基态的特征是：粒子对的波函数在大于 (11.73) 所定义的相干长度 $\xi$ 的距离上存在强关联。因此，让我们回到 §11.7 的 London 理论，并且有意为电子密度的超流分量引入一个\emph{宏观波函数} $\Psi(\mathbf{r})$（通常称为 Ginzburg--Landau 序参量（order parameter））。正如我们在 §11.8 中学到的那样，只要我们不涉及在短于相干长度的距离上迅速变化的性质，这可以很好地充当超流序的一个适当的局域量度。

按照量子力学的正则形式，我们自然假设局域超流密度由
\begin{equation*}
n_s(\mathbf{r}) = |\Psi(\mathbf{r})|^{2} \tag{11.79}
\end{equation*}
给出，并且
\begin{equation*}
\mathbf{J}_s(\mathbf{r}) = \frac{e^{*}}{2m^{*}}\left[\Psi^{*}\left\{\frac{\hbar}{i}\nabla - \frac{e^{*}}{c}\mathbf{A}\right\}\Psi + \Psi\left\{-\frac{\hbar}{i}\nabla - \frac{e^{*}}{c}\mathbf{A}\right\}\Psi^{*}\right] \tag{11.80}
\end{equation*}
定义超流电流密度，正如 (10.3) 和 (11.55) 那样。但由于怀疑超流电流的载体是“对”而不是单个电子，我们让 $e^{*}$ 和 $m^{*}$ 保持未定义。

如果 $\Psi(\mathbf{r})$ 处处为实，那么它就可以被消去而换成 $n_s$：(11.79) 和 (11.80) 立即约化为 London 方程 (11.57)，只是把 $e$、$m$ 和 $n$ “重整化”为 $e^{*}$、$m^{*}$ 和 $n_s$。在 §11.8 所讨论的限制之内，这就为 Meissner 效应提供了一个良好的定性描述，其中包括隐含在 $n_s$ 随 $T$ 变化之中的穿透深度对温度的依赖。

但是，除非 $\Psi(\mathbf{r})$ 是\emph{复}的，我们就不能描述均匀材料中的超流电流。这一唯象理论出乎意料的特征是：物理情况直接依赖于序参量的\emph{相位}的变化，而不只是依赖于它的绝对大小：用矩阵力学的语言说，我们必须与非对角长程序（off-diagonal long range order）打交道。

例如，假设我们有一种温度恒定且均匀的材料，其中 $n_s(\mathbf{r})$ 为常数。写成
\begin{equation*}
\Psi(\mathbf{r}) = (n_s)^{\frac{1}{2}}\,e^{i\chi(\mathbf{r})}, \tag{11.81}
\end{equation*}
我们便得到
\begin{equation*}
m^{*}\mathbf{v}_s(\mathbf{r}) = \hbar\nabla\chi(\mathbf{r}) - (e^{*}/c)\,\mathbf{A}(\mathbf{r}) \tag{11.82}
\end{equation*}
在不存在磁场时，相位函数 $\chi(\mathbf{r})$ 扮演 (11.77) 中已定义的超流速度 $\mathbf{v}_s$ 的速度势的角色。但如果 $\chi$ 不随空间变化，就不会有超流电流。

\begin{figure}[H]
\centering
\includegraphics[width=0.72\textwidth]{fig_209.pdf}
\caption*{图 209\quad 用于磁通量子化（flux quantization）的积分路径。}
\end{figure}

现在考虑一个被磁场（具有矢势 $\mathbf{A}(\mathbf{r})$）穿过的超导材料回路（图 209）。让我们沿环绕此电路的一条路径积分 (11.82)，路径保持在超导介质内部足够深处，以避开可能发生抗磁屏蔽电流的穿透区。在电路中不存在其他电动势源的情况下，$\mathbf{v}_s(\mathbf{r})$ 沿此
路径必须为零。序参量 $\Psi(\mathbf{r})$ 本身必须是超导体内位置的单值函数，因此当我们绕回到出发点时，其相位角必须回到同一数值，再加上 $2\pi$ 的某个整数倍。由初等矢量演算可得
\begin{equation*}
\Phi = \oint \mathbf{A}\cdot\mathbf{dl} = (\hbar c/e^{*}) \oint \nabla\chi(\mathbf{r})\cdot\mathbf{dl} = n\,2\pi\hbar c/e^{*},
\tag{11.83}
\end{equation*}
其中 $n$ 为整数。\emph{全磁通}（fluxoid）$\Phi$——即穿过回路的总磁通——以 $2\pi\hbar c/e^{*}$ 为单位而量子化。

\emph{磁通量子化}（flux quantization）与液 $^{4}$He 中类似的\emph{环流量子化}（quantization of circulation）一道，是量子力学最引人注目的宏观表现之一。实验发现，冻结在超导系统中的磁通以下列量子为单位跳变
\begin{equation*}
\Phi_0 = 2\pi\hbar c/2e = 2.07\times 10^{-7}\,\mathrm{G\,cm^{2}}.
\tag{11.84}
\end{equation*}
超流粒子携带的电荷恰好是一个电子对的电荷，即
\begin{equation*}
e^{*} = 2e.
\tag{11.85}
\end{equation*}
除了这一重整化之外，基本的 London 假设由此得到完全的实验证实。

对电荷 (11.85) 的确正是我们根据 BCS 理论所应预期的结果。要从第一性原理推导诸如 (11.80) 那样的方程（同时顾及“局域性”近似），需要依靠 \emph{Gor'kov 方程}，它以 Green 函数的语言写成。在这一形式体系中，序参量 $\Psi(\mathbf{r})$ 联系着一个“反常传播子”（anomalous propagator）
\begin{equation*}
F_{\uparrow\downarrow}(\mathbf{r}, t; \mathbf{r}', t') = \langle \psi_\uparrow(\mathbf{r}, t)\,\psi_\downarrow(\mathbf{r}', t') \rangle
\tag{11.86}
\end{equation*}
它度量的是自旋相反的一对电子的波函数在 $\mathbf{r}$ 与 $\mathbf{r}'$ 两点、$t$ 与 $t'$ 两时刻之间的平均相位相干性。正如我们在 (11.76) 中看到的，这个量不消失正是超流态的特征；而在金属的“正常”态中，它本来会自动消失。

\section{超导结}

超导体中的能隙可以通过 Giaever 隧穿（Giaever tunnelling）直接观测。如 §6.8 中那样，我们研究夹在譬如正常金属与超导体之间的一层薄绝缘膜的电流/电压特性。正常电子只有当有空着的准粒子态来接纳它们时才能进入超导体。在远低于 $T_c$ 的温度下，电流必须为零，直到外加偏压 $\phi$ 达到能隙 $\Delta$（§11.4），然后急剧跳升（图 210）。这种结的透射系数 $\mathcal{T}$ 难以计算，但电流随电压的变化使我们得以细致检验能隙区中准粒子态的态密度。

\begin{figure}[H]
\centering
\includegraphics[width=0.72\textwidth]{fig_210.pdf}
\caption*{图 210\quad 准粒子从正常金属隧穿进入超导体：(a) 零偏压时；(b) $\phi > \Delta$ 时；(c) 电流/电压特性，显示温度的影响。}
\end{figure}

实际上，由于 §6.8 中已经讨论过的原因，这一转变并不完全陡峭。对于强耦合超导体，即 (11.22) 中参量 $\mathcal{N}(\mathcal{E}_F)V$ 很大的情形，电子–声子相互作用 (11.5) 足以在准粒子中产生相当比例的\emph{声子协助隧穿}（phonon-assisted tunnelling）（参看图 116）。于是，在隧穿电流中观察到的精细结构可以给出关于超导体声子谱的详细信息。

更引人注目的现象是\emph{相干隧穿}（coherent tunnelling），此时有超流电流流过结。它不仅可以通过两个超导体之间的薄绝缘膜来观测，还可以透过正常金属中相当厚的一层（5000 Å）来观测——借助邻近效应（proximity effect），这层正常金属可以暂时变成超导的。换句话说，真正超导体中配对相互作用的强度，足以在电子穿过势垒时保持其间的相位关联，并在另一侧重新建立起相干性。

对于所有这类现象，§11.9 的宏观序参量形式体系显然是合用的。这个函数在这样的边界上如何行为？结左方的 $\Psi_L$ 与右方的 $\Psi_R$ 之间有什么关系？

既然穿过势垒的渗漏是一个随时间变化的过程，我们就需要考虑 $\Psi$ 的运动方程。按正则量子力学的说法，这涉及一个超流粒子的能量，我们把它取为某个值 $\mu$。于是我们为序参量写出一个含时 Schrödinger 方程，即
\begin{equation*}
i\hbar\frac{\partial\Psi}{\partial t} = \mu\Psi.
\tag{11.87}
\end{equation*}
当然，在普通的超流电流中，$\mu$ 应当是超流配对的化学势——即在 (11.27) 中用作能量原点的量 $\zeta$ 的两倍——因而在整个回路上严格处处相同。这样的系统因此在时间上表现为定态，正如 (11.80) 诸式所隐含的那样。

然而，对于通过两个相同超导体之间的结的透射，我们必须考虑耦合方程
\begin{equation*}
\left.\begin{array}{l}
i\hbar\dfrac{\partial\Psi_L}{\partial t} = \mu_L\Psi_L + \mathcal{T}\Psi_R,\\[8pt]
i\hbar\dfrac{\partial\Psi_R}{\partial t} = \mu_R\Psi_R + \mathcal{T}\Psi_L,
\end{array}\right\}
\tag{11.88}
\end{equation*}
其中我们容许超流序以速率 $\mathcal{T}$ 双向转移。注意，这个系数在这里以\emph{矩阵元}的身份出现，虽然它起初定义为准粒子隧穿的\emph{跃迁概率}。这是因为超流由\emph{电子对}组成，电子对的隧穿速率预期正比于单电子速率的\emph{平方}。

方程 (11.88) 具有一般的含时解，其中超流密度 $|\Psi_L|^{2}$ 与 $|\Psi_R|^{2}$ 以相等而相反的速率变化，仿佛超流粒子正从结的一侧被输运到另一侧。我们应当把这个解解释为一个（单位面积的）电流密度
\begin{equation*}
J_s = e^{*}\frac{\partial}{\partial t}|\Psi_R|^{2} = \frac{2}{\hbar}\,\mathcal{T}\,n_s e^{*}\sin(\chi_R - \chi_L),
\tag{11.89}
\end{equation*}
其中 $\Psi_R$ 与 $\Psi_L$ 之间的相位差（参看 (11.81)）随时间按下列方式变化
\begin{equation*}
\frac{\partial}{\partial t}(\chi_R - \chi_L) = -\frac{1}{\hbar}(\mu_R - \mu_L).
\tag{11.90}
\end{equation*}
当然，这些方程可以在 Gor'kov 形式体系中从 BCS 理论严格推导出来，只需用一个隧穿哈密顿量来表示电子从一个区域向另一个区域的转移。这里注意，超流电流 (11.89) 与准粒子隧穿系数 $\mathcal{T}$ 成正比，这表明相位相干性再一次发挥了作用。

现在假设我们在结两端加上一个小的恒定电压 $\phi$（小于能隙宽度，以避免准粒子激发）。这相当于超流粒子的化学势有了 $2e\phi$ 的差。方程 (11.89) 与 (11.90) 描写一个随时间振荡的超流电流，其频率为 $2e\phi/\hbar$。这就是交流 Josephson 效应（A.C. Josephson effect），它再次为量子现象中的相位相干性提供了直接的宏观证据，并且已被证明是测量重要原子常数之比 $e/\hbar$ 的极为精确的方法。

在磁场存在时，这些公式需要按式 (10.3) 作通常的修改，以保持矢量势定义中的规范不变性。一个典型的情形如下。假设两个超导体之间有一个面积较大的结（图 211）。由于绝缘势垒本身不是超导的，磁通可以穿过结区，方向平行于分隔面。今考虑成对的点，如 $L_1,R_1$；$L_2,R_2$，位于结的两侧，深入到各自内部足以避开穿透区，但在平行于表面的方向上相距颇远。假设已知相位差恰为
\begin{equation*}
\Delta\chi_1 = \chi(R_1) - \chi(L_1)
\tag{11.91}
\end{equation*}
——它跨接 $L_1R_1$；那么跨接 $L_2R_2$ 的相应相位差 $\Delta\chi_2$ 是多少？

考虑一条从 $L_2$ 到 $L_1$ 的路径，沿这条路径我们对式 (11.82) 的两边积分。正如推导 (11.83) 时那样，$\mathbf{v}_s$ 没有贡献——除穿透区之外，它不可能有平行于结面的分量：于是得到
\begin{equation*}
\chi(L_1) - \chi(L_2) = \frac{e^{*}}{\hbar c}\int_{L_2}^{L_1}\mathbf{A}(\mathbf{r})\cdot\mathbf{dl}.
\tag{11.92}
\end{equation*}

\begin{figure}[H]
\centering
\includegraphics[width=0.72\textwidth]{fig_211.pdf}
\caption*{图 211\quad Josephson 结。}
\end{figure}

类似地，在右边的材料中沿返回的路径，
\begin{equation*}
\chi(R_2) - \chi(R_1) = \frac{e^{*}}{\hbar c}\int_{R_1}^{R_2}\mathbf{A}(\mathbf{r})\cdot\mathbf{dl}.
\tag{11.93}
\end{equation*}
把这两个方程相加，并略去穿过结时右方路径上的平庸贡献，我们得到
\begin{align*}
\Delta\chi_2 &= \Delta\chi_1 + \frac{e^{*}}{\hbar c}\oint \mathbf{A}\cdot\mathbf{dl}\\
&= \Delta\chi_1 + 2\pi\Phi(1,2)/\Phi_0,
\tag{11.94}
\end{align*}
其中 $\Phi_0$ 是量子化磁通的单位，即 (11.84)，而 $\Phi(1,2)$ 是穿过势垒、介于直线 $L_1R_1$ 与直线 $L_2R_2$ 之间的磁通。换言之，当我们在结区内四处移动时，序参量跨边界的相对相位差是变化的。

于是，宽结上的总电流是外加磁场的复杂函数。我们必须对每一小块面积计算 $(\chi_R - \chi_L)$，由 (11.89) 算出对电流的贡献，然后再对整个结面求积分。例如，对于宽度为 $W$、被总磁通 $\Phi$ 穿过的结，可以导出最大 Josephson 电流的公式，形式为
\begin{equation*}
J \propto \left|\sin(\pi\Phi/\Phi_0)\big/(\pi\Phi/\Phi_0)\right|,
\tag{11.95}
\end{equation*}
这一电流可以在零外加电压下流动，但当改变磁场时它迅速振荡。这种直流 Josephson 效应（D.C. Josephson effect）在双结并联这一更复杂的构型中，提供了测量磁场的极其灵敏的器件。

\section{第二类材料}

对 Meissner 效应理论（§11.7）作细致的分析会揭示一个佯谬。我们曾假设，只要磁场能量一超过超导态与正常态的自由能之差，整块样品就会停止超导——亦即在式 (11.63) 给出的临界场 $H_c$ 处。但这个假设忽略了构造一种非均匀相的可能性：磁通由弥散在超导材料中的正常金属细丝层来承载。这样的\emph{中间态}（intermediate state）显然直到高得多的磁场仍可保持热力学稳定，并且在宏观尺度上仍然是电学超导的。

实验上，在\emph{第二类材料}（type II materials）中观察到的正是与此相仿的状态：Meissner 效应在较低的下临界场 $H_{c1}$ 处消失，但材料直到高得多的场 $H_{c2}$ 仍不显示电阻。这类材料用作电磁铁绕组的重要性使人们的注意力集中于这一现象；这一现象在理论上也有重大的意义。

正如铁磁畴理论（§10.3）中那样，我们需要为同一金属中正常区与超导区之间的边界建立一个理论模型。这意味着，(11.77) 或 (11.79) 中的“超流密度”$n_s$ 必须允许随空间变化。我们还需要一些进一步的方程，来确定序参量 $\Psi(\mathbf{r})$ 的振幅作为 $\mathbf{r}$ 的函数的行为。

Ginzburg--Landau 理论的基本假设是：处于磁场 $H$ 中的超导体，其自由能密度（相对于同温度下零场中的正常金属来计量）是如下形式的函数
\begin{equation*}
g_s - g_n = \alpha|\Psi|^{2} + \tfrac{1}{2}\beta|\Psi|^{4} + \frac{1}{8\pi}H^{2} + \frac{1}{2m^{*}}\left|-i\hbar\nabla\Psi - \frac{e^{*}}{c}\mathbf{A}\Psi\right|^{2}.
\tag{11.96}
\end{equation*}
头两项仅仅表达了一个任意的假设：进入超导相所获得的自由能收益是序参量 $\Psi$ 的函数，并且可以按实量 $|\Psi|^{2}$ 的幂展开。任意系数 $\alpha$ 与 $\beta$ 本身依赖于温度；在 $T_c$ 以下、超导相占优的温度范围内，系数 $\alpha$ 必然变为负；于是零磁场中的超流浓度受 $\beta$ 项的限制，取值为
\begin{equation*}
n_s = |\Psi_0|^{2} = -\alpha/\beta.
\tag{11.97}
\end{equation*}
对于遵从 Meissner 效应的第一类材料，我们可以把 $H^{2}$ 项考虑进来，并把大块材料的临界场认定为
\begin{equation*}
H_c^{2}/8\pi = \alpha^{2}/2\beta,
\tag{11.98}
\end{equation*}
亦即 $g_s$ 再也无法降到 $g_n$ 之下的那一点。同样，我们可以把 (11.97) 代入 London 公式 (11.61)，即
\begin{equation*}
\lambda = \left(\frac{m^{*}c^{2}}{4\pi n_s e^{*2}}\right)^{\frac{1}{2}},
\tag{11.99}
\end{equation*}
以便用同一组参量来表达穿透深度，从而可以估计这些参量的温度依赖性。

(11.96) 的最后一项也是唯象的，但它有牢固的正则基础——它正是磁场中粒子动能的典型量子力学表达式。在 London 方程 (11.55) 与 (11.80) 中，这一项代表“超流电流”的动能，它来自 $\Psi$ 相位的空间变化，如 (11.82) 那样。Ginzburg--Landau 理论的精妙之处在于：恰恰是同一项，系数不改，却能同时描述序参量的模发生空间变化所产生的效应——这些效应本应归因于破坏配对相干性所需的能量。

例如，假设我们想从一个超导区——其中 $\Psi$ 由 (11.97) 给出——过渡到一个 $\Psi = 0$ 的正常区，同时使总自由能的代价最小。在没有磁场并且忽略 $\beta$ 项的情况下，序参量的模必须满足偏微分方程
\begin{equation*}
(\hbar^{2}/2m^{*})\nabla^{2}|\Psi| - |\alpha|\Psi = 0,
\tag{11.100}
\end{equation*}
它由自由能积分经变分法得到。这描写的是 $|\Psi|$ 指数式衰减的情形，其特征长度为
\begin{equation*}
\xi = \left\{\hbar^{2}/2m^{*}|\alpha|\right\}^{\frac{1}{2}}.
\tag{11.101}
\end{equation*}
换言之，我们可以把系数 $\alpha$ 与 §11.8 的相干长度联系起来；在典型的“脏”金属中，这个量近似等于电子平均自由程 $\Lambda$。这个描述显然不足以细致地表现 §11.8 中 Pippard 公式所隐含的非局域行为，但在接近 $T_c$ 的温度下，可以用 Gor'kov 方法从第一性原理予以论证。

论证的下一步，是研究超导材料（其中 $|\Psi| = |\Psi_0|$ 且 $H = 0$）与含有临界场 $H_c$ 的正常材料之间的临界边界处自由能最小的位形。$\Psi$ 的一般偏微分方程要与 $H$ 的 Maxwell 方程组相耦合，如 (11.58)–(11.60) 那样，是相当繁复的，但在两种极端情形下，答案的形式不难求出。

首先假设相干长度 $\xi$ 远大于穿透深度 $\lambda$。由图 212($a$) 可以清楚看到，磁场被从厚度为 $\xi$ 的一层中排除出去，这需要能量
\begin{equation*}
\sigma_{\mathrm{I}} \sim \xi(H_c^{2}/8\pi)
\tag{11.102}
\end{equation*}
（每单位边界面积），而在这一层的大部分区域内，并没有“有序化”自由能的收益作为补偿。表面能近似由 (11.102) 给出，是正的。这是第一类超导电性的典型情形：中间态因需要在正常区与超导区之间制造边界而付出能量，故受到抑制。

另一方面，如果相干长度 $\xi$ 远小于穿透深度 $\lambda$，那么在磁场被排除的区域的大部分范围内，都可以获得有序相的负能量（图 212($b$)）。因此粗略地说，表面能的形式为
\begin{equation*}
\sigma_{\mathrm{II}} \sim -\lambda(H_c^{2}/8\pi).
\tag{11.103}
\end{equation*}
更细致的分析表明，表面能在下述情况下由正变负：
\begin{equation*}
\lambda > \xi/\sqrt{2}.
\tag{11.104}
\end{equation*}
这就是第二类超导电性的基本条件；化学杂质与晶体缺陷显然有利于这一条件，因为它们使相干长度 $\xi$ 缩短，而对 $\lambda$ 等其他宏观超导参量没有明显的影响。

\begin{figure}[H]
\centering
\includegraphics[width=0.72\textwidth]{fig_212.pdf}
\caption*{图 212\quad 正常区与超导区之间的边界。(a) 第一类，$\xi \gg \lambda$；(b) 第二类，$\xi \ll \lambda$。}
\end{figure}

正常相与超导相之间的负表面能显然有利于第二类材料在 $H_{c1}$ 与 $H_{c2}$ 之间的磁场范围内所采取的\emph{混合态}（mixed state）。样品被磁通细丝所穿透，每一根细丝都具有单位全磁通强度，即式 (11.84)。细丝内有一个半径为 $\xi$ 的正常金属圆柱形芯区，外围是半径为 $\lambda$ 的圆柱区域，超流电流必须围绕它流动，以限制磁场的穿透。因此，每根细丝的行为都像超流体中的一根涡线（vortex line），并且倾向于排斥其近邻。当我们增大磁场时，磁通线的密度增大，细丝“结晶”成二维点阵。但间隙材料在宏观尺度上仍保持超导，直到上临界场 $H_{c2}$——此时正常金属的芯区被迫相互接触，整块样品回归正常态。

\begin{figure}[H]
\centering
\includegraphics[width=0.72\textwidth]{fig_213.pdf}
\caption*{图 213\quad 第二类超导体中磁通细丝的阵列，显示磁力线集中于圆柱形正常材料芯的内部及其周围。超流涡旋在穿透区内绕细丝流动。}
\end{figure}
